\documentclass[11pt]{article}
% Adaptado de C:/GitHub/Mestrado/Template/Template.tex.
\usepackage[T1]{fontenc}
\usepackage[utf8]{inputenc}
\usepackage[brazilian]{babel}
\usepackage{graphicx,booktabs,array,amsmath,amssymb,microtype}
\usepackage{newtxtext,newtxmath}
\usepackage{geometry}
\usepackage{pdflscape,caption}
% Apresentação de artigo para leitura em tela, conforme preferência do autor.
\geometry{a4paper,margin=2.5cm,headheight=12pt,headsep=0.5cm}
\usepackage{setspace,indentfirst,csquotes}
\setstretch{1.05}
% Parágrafos compactos, com recuo discreto e sem espaço adicional.
\setlength{\parindent}{0.5cm}
\setlength{\parskip}{0pt}
\usepackage{titlesec}
\titleformat{\section}[hang]{\normalfont\large\bfseries}{\thesection}{\fontdimen2\font}{}
\titleformat{name=\section,numberless}[block]{\normalfont\large\bfseries\filcenter}{}{0pt}{}
\titleformat{\subsection}[hang]{\normalfont\normalsize\bfseries\itshape}{\thesubsection}{\fontdimen2\font}{}
\titlespacing{\section}{0pt}{0pt}{8pt}
\titlespacing{\subsection}{0pt}{10pt}{5pt}
\titlespacing{\paragraph}{0pt}{7pt}{0.7em}
\AtBeginDocument{%
  \setlength{\abovedisplayskip}{8pt plus 2pt minus 2pt}%
  \setlength{\belowdisplayskip}{8pt plus 2pt minus 2pt}%
  \setlength{\abovedisplayshortskip}{4pt plus 2pt}%
  \setlength{\belowdisplayshortskip}{6pt plus 2pt minus 2pt}%
}
\usepackage{fancyhdr}
\pagestyle{fancy}
\fancyhf{}
\fancyhead[R]{\footnotesize\thepage}
\renewcommand{\headrulewidth}{0pt}
\renewcommand{\footnoterule}{\kern-3pt\hrule width 5cm\kern2.6pt}
\makeatletter
\long\def\@makefntext#1{\parindent 0pt\noindent\hangindent=1.8em\hangafter=1\hb@xt@1.8em{\hss\@makefnmark}#1}
\makeatother
% Notas e legendas em tamanho 9; espaço simples nos objetos.
\DeclareCaptionFont{simples}{\setstretch{1}}
\captionsetup{font={footnotesize,simples},labelsep=endash,justification=raggedright,singlelinecheck=false,position=top,skip=6pt}
\renewenvironment{abstract}{\par\medskip\noindent\textbf{\abstractname}\par\smallskip\small\setstretch{1}\noindent}{\par\medskip}
\interfootnotelinepenalty=10000
\widowpenalty=10000
\clubpenalty=10000
\displaywidowpenalty=10000
\usepackage[backend=biber,style=abnt]{biblatex}
\addbibresource{references.bib}
\setlength{\bibhang}{0pt}
\renewcommand*{\bibfont}{\normalfont\small\setstretch{1}}
\AtBeginBibliography{\setlength{\bibitemsep}{5pt}}
\DeclareFieldFormat{url}{\small\bibstring{urlfrom}\addcolon\addspace\url{#1}}
\DeclareFieldFormat{urldate}{\small\bibstring{urlseen}\addcolon\addspace #1}
\usepackage{hyperref}
\hypersetup{hidelinks}
\newcommand{\instituicao}{UNIVERSIDADE FEDERAL DE SANTA CATARINA}
\newcommand{\departamento}{Programa de Pós-Graduação em Economia}
\newcommand{\curso}{Mestrado}
\newcommand{\disciplina}{Desigualdade, Diversidade e Redes}
\newcommand{\titulo}{Rede de emissões de CO\textsubscript{2} no Brasil}
\newcommand{\autor}{Gabriel Pincelli Cintra}
\newcommand{\local}{Florianópolis}
\newcommand{\dataentrega}{18 de outubro de 2026}
\newcommand{\imprimirtitulodois}{
\begin{center}
  {\small\MakeUppercase{\instituicao}\par}
  {\small \departamento\par}
  {\small \disciplina\ (\curso)\par}
  \vspace{1em}{\LARGE\bfseries \titulo\par}
  \vspace{0.7em}{\normalsize \autor\par}
  \vspace{0.3em}{\small \local, \dataentrega\par}
\end{center}
\vspace{0.4em}}
\begin{document}
\imprimirtitulodois
\thispagestyle{empty}
\begin{abstract}
Este ensaio examina a distribuição das emissões de CO\textsubscript{2} nas relações entre atividades emissoras e setores de produtos finais no Brasil. Combina a Matriz de Insumo-Produto de 2015, com 67 atividades, e intensidades de CO\textsubscript{2} em uma rede dirigida e ponderada. As forças e o número efetivo de contrapartes caracterizam volume e concentração; o filtro de disparidade seleciona ligações conforme seu peso relativo nas emissões de cada origem ou destino. Apesar da conectividade quase uniforme, energia e água/resíduos distribuem volumes externos próximos com concentrações distintas, enquanto construção reúne emissões de diversas origens. As contrapartes esclarecem esses perfis e a importância assimétrica de uma mesma relação para origem e destino. Na construção, 39,5\% da conta pela demanda corresponde a emissões geradas em outras atividades. O backbone, sub-rede formada pelas ligações selecionadas, retém 13,40\% das ligações e 70,69\% do peso intersetorial, com coberturas locais desiguais. A análise oferece elementos para discutir eficiência material e aquisição de materiais de menor emissão na cadeia da construção, com interpretações condicionadas à validação dos coeficientes.
\end{abstract}
\noindent\textbf{Palavras-chave:} atribuição de emissões; insumo-produto; redes dirigidas e ponderadas; mitigação.

\clearpage
\input{introducao.tex}
% Sequência inspirada em Huang et al. (2019).
\clearpage
\input{dados.tex}
\clearpage
\section{Metodologia}
\input{modelo_atribuicao.tex}
\subsection{Construção da rede dirigida e ponderada}
\label{sec:construcao-rede}

A rede tem 67 nós, um por atividade da MIP, identificados pelos códigos do IBGE. Cada nó representa uma origem emissora nas ligações de saída e o setor de um produto final nas ligações de entrada, conforme a representação de \textcite{huang2019}. Em cada nó de destino, a atribuição reúne as emissões associadas aos diferentes usos finais dos produtos da atividade, como consumo das famílias e exportações.

Para estudar relações entre atividades distintas, define-se
\begin{equation}
W_{ij}=\begin{cases}
F_{ij}, & i\ne j,\\
0, & i=j.
\end{cases}
\label{eq:rede-sem-diagonal}
\end{equation}
Cada $W_{ij}>0$ define uma aresta $i\rightarrow j$, cujo peso, em Gg de CO\textsubscript{2}, indica a emissão em $i$ atribuída ao produto final de $j$. As duas direções de um par são distintas. Como os pesos incorporam requisitos diretos e indiretos, concatenar arestas não representa uma sequência de etapas produtivas.

A diagonal $F_{ii}$ reúne emissões de $i$ atribuídas à demanda final de seus próprios produtos, incluindo requisitos indiretos. Ela permanece em $F$ e nas contas $c$ e $u$, mas é retirada de $W$, cujo peso total corresponde apenas às atribuições entre atividades distintas.

A construção preserva todos os nós e as arestas positivas, sem corte por magnitude. São conferidos o alinhamento dos códigos e a existência de pesos finitos e não negativos. A seleção de arestas é tratada na seção~\ref{sec:backbone}.

\subsection{Volume e concentração das relações}
\label{sec:volume-concentracao}

O volume e a concentração são calculados na rede completa $W$, por meio das forças e da disparidade dos pesos \parencite{huang2019,serrano2009}.

\paragraph{Volume.}
As forças de saída e entrada são as somas dos pesos em cada direção:
\begin{equation}
\begin{aligned}
s_i^{\mathrm{out}}&=\sum_{j\ne i}W_{ij}=c_i-F_{ii},\\
s_i^{\mathrm{in}}&=\sum_{j\ne i}W_{ji}=u_i-F_{ii}.
\end{aligned}
\label{eq:forcas-externas}
\end{equation}
A força de saída mede as emissões próprias de $i$ atribuídas aos produtos finais de outras atividades; a de entrada mede as emissões de outras atividades atribuídas aos produtos finais de $i$. A diferença entre força de entrada e força de saída é $s_i^{\mathrm{in}}-s_i^{\mathrm{out}}=u_i-c_i$. A leitura separada dos volumes revela atribuições que se compensam no saldo; as contas completas e a diagonal situam sua participação no total.

\paragraph{Concentração.}
Para perfis com força positiva, define-se a participação de cada contraparte:
\begin{equation}
p_{ij}^{\mathrm{out}}=\frac{W_{ij}}{s_i^{\mathrm{out}}},\qquad
p_{ji}^{\mathrm{in}}=\frac{W_{ji}}{s_i^{\mathrm{in}}},\quad j\ne i.
\label{eq:participacoes-externas}
\end{equation}
O perfil de saída distribui o volume intersetorial da origem entre destinos; o de entrada descreve as origens que compõem a atribuição intersetorial ao destino. Cada perfil soma um. A disparidade e seu inverso são
\begin{equation}
\begin{aligned}
Y_i^{\mathrm{out}}&=\sum_{j\ne i}(p_{ij}^{\mathrm{out}})^2,&
Y_i^{\mathrm{in}}&=\sum_{j\ne i}(p_{ji}^{\mathrm{in}})^2,\\
D_{2,i}^{\mathrm{out}}&=1/Y_i^{\mathrm{out}},&
D_{2,i}^{\mathrm{in}}&=1/Y_i^{\mathrm{in}}.
\end{aligned}
\label{eq:disparidade}
\end{equation}
$Y$ aumenta com a concentração; $D_2$ expressa a mesma informação como número efetivo de contrapartes. Para $k$ pesos iguais, $D_2=k$; quando uma relação domina, aproxima-se de um. As medidas são adimensionais e ficam indefinidas quando a força é zero, sem excluir o nó.

Os pesos e as participações identificam as contrapartes, cuja identidade não é informada por $D_2$. O grau conta as arestas positivas de entrada ou saída; a densidade é o número de arestas dividido por $67\times66$. Esses indicadores descrevem a presença de relações, enquanto a disparidade considera a distribuição dos pesos.

\subsection{Extração e avaliação do backbone}
\label{sec:backbone}

O filtro de disparidade é um método que reduz a rede ao selecionar ligações com participações excepcionalmente altas nos pesos de uma origem ou de um destino, segundo um modelo nulo \parencite{serrano2009,huang2019}. A sub-rede formada pelas ligações selecionadas é denominada backbone e preserva seus pesos originais, sem alterar as contas ambientais completas.

\paragraph{Critério de seleção.}
O modelo nulo distribui a força entre $k$ ligações dividindo o intervalo unitário por $k-1$ pontos uniformemente escolhidos. Para $k>1$, a probabilidade de uma parcela atingir ou superar $p$ é $(1-p)^{k-1}$. Para a aresta $i\rightarrow j$, as caudas nas perspectivas de saída e entrada são
\begin{equation}
\begin{aligned}
a_{ij}^{\mathrm{out}}&=\left(1-\frac{W_{ij}}{s_i^{\mathrm{out}}}\right)^{k_i^{\mathrm{out}}-1},\\
a_{ij}^{\mathrm{in}}&=\left(1-\frac{W_{ij}}{s_j^{\mathrm{in}}}\right)^{k_j^{\mathrm{in}}-1}.
\end{aligned}
\label{eq:caudas-filtro}
\end{equation}
Os graus $k$ e as forças são os de $W$. A aresta é selecionada se $a_{ij}^{\mathrm{out}}<\alpha$ \textbf{ou} $a_{ij}^{\mathrm{in}}<\alpha$, registrando-se a extremidade responsável pela seleção. A saída destaca um destino nas atribuições da origem; a entrada destaca uma origem na composição do destino.\footnote{As caudas são calculadas apenas para graus pertinentes maiores que um. Se uma extremidade tem grau um, vale o teste da outra; se ambas têm grau um, a aresta é preservada por convenção de manutenção da relação exclusiva, inclusive em díades isoladas. Essa extensão dos casos degenerados do suplemento de Serrano et al. é registrada separadamente da seleção estatística.}

Adota-se $\alpha=0{,}05$, com $0{,}01$ e $0{,}10$ para sensibilidade. O filtro é aplicado uma vez a $W$ para cada valor, sem atualização dos perfis. A matriz $B_\alpha$ preserva os pesos originais nas arestas selecionadas e zero nas demais, sem renormalização ou corte absoluto adicional.

\paragraph{Avaliação da seleção.}
Comparam-se a fração de arestas retidas e a cobertura dos pesos, globalmente e por atividade:
\begin{equation}
\begin{aligned}
r_\alpha&=\frac{\sum_{ij}(B_\alpha)_{ij}}{\sum_{ij}W_{ij}},\\
r_{\alpha,i}^{\mathrm{out}}&=\frac{\sum_j(B_\alpha)_{ij}}{s_i^{\mathrm{out}}},\qquad
r_{\alpha,i}^{\mathrm{in}}=\frac{\sum_j(B_\alpha)_{ji}}{s_i^{\mathrm{in}}}.
\end{aligned}
\label{eq:cobertura-backbone}
\end{equation}
As razões são definidas apenas para denominadores positivos. A cobertura refere-se à parcela intersetorial de $W$ e permite distinguir retenção agregada de perdas por atividade e direção. A conectividade é descrita pelos graus, isolados e componentes fracamente conectados, obtidos ao desconsiderar a direção das arestas.

A sensibilidade é examinada pelas mudanças nas arestas, coberturas e posições setoriais entre os três valores de $\alpha$. O aninhamento dos conjuntos decorre da regra e não constitui evidência independente de robustez. As caudas do modelo nulo não medem incerteza ambiental, e a seleção não estabelece prioridade de mitigação.

Para interpretar os perfis, as contrapartes são ordenadas por peso em $W$ completa. Destacam-se as cinco maiores e a participação conjunta das demais, sem modificar a rede. A leitura dos casos combina essas participações, as contas completas e os motivos de retenção; a sensibilidade verifica também a permanência das principais relações nos três parâmetros.

\clearpage
\input{resultados.tex}
\clearpage
\section{Discussão e implicações para a coordenação da mitigação}

A diferença entre as contas pela produção e pela demanda informa o saldo de uma atividade, mas não identifica as contrapartes nem a concentração dos volumes que o compõem. Energia e água/resíduos atribuem a outros produtos finais volumes superiores aos provenientes de outras origens; em construção, predomina o sentido inverso. A leitura conjunta das contas e dos perfis esclarece essas diferenças: a diagonal mantém a parcela atribuída aos produtos finais da própria atividade, e as relações intersetoriais identificam as emissões atribuídas a setores distintos de sua origem. Essa leitura complementa as comparações de atribuição entre setores, regiões e cadeias produtivas presentes na literatura brasileira \parencite{imori2015,sanguinetAzzoni2024,montoyaEtAl2026}.

Volumes semelhantes também podem envolver conjuntos de contrapartes com participações diferentes. A maior concentração das saídas de água/resíduos, em comparação com energia, destaca o peso da administração pública no primeiro perfil. Já a concentração das saídas de construção não caracteriza suas entradas, mais distribuídas entre origens. Esses contrastes mostram por que a posição de uma atividade nas contas não basta para descrever sua participação nas relações intersetoriais. A representação dirigida de \textcite{huang2019} permite distinguir os dois papéis de cada atividade, enquanto a análise dos pesos considera diferenças que o grau quase uniforme da rede completa não informa.

A importância local distingue, ainda, relações destacadas para a origem ou para o destino. Construção pode representar uma parcela expressiva das atribuições externas de produtos de metal sem que essa origem domine a composição recebida pela construção. Assim, examinar apenas os maiores componentes da conta de um destino pode deixar em segundo plano relações importantes na perspectiva das origens. O filtro de \textcite{serrano2009} explicita essa assimetria sob seu modelo nulo. A seleção precisa ser lida com as participações e as coberturas, pois as ligações omitidas podem representar, em conjunto, parcela expressiva de um perfil.

Em uma aplicação a Santa Catarina em 2018, esse procedimento permitiria complementar a comparação das participações setoriais com duas questões: quais produtos finais estão associados às emissões dos setores cuja participação diminui na atribuição pela demanda e quais origens compõem a atribuição dos setores cuja participação aumenta. A concentração distinguiria perfis com volumes mais equilibrados entre contrapartes daqueles dominados por poucas relações. A aplicação exigirá dados harmonizados e o recálculo das contas e dos perfis na fronteira estadual; as contrapartes e os contrastes do cenário nacional não podem ser transportados para o estado.

A parcela intersetorial identificada na construção justifica considerar as emissões geradas na produção dos insumos ao avaliar a mitigação associada aos produtos finais desse setor. Entre as origens identificadas estão produtos de minerais não metálicos, energia e produtos de metal. O diagnóstico pode ser relacionado às estratégias de eficiência material e aos critérios de emissões incorporadas nos projetos de construção \parencite[seções~5.1 e~6.4]{rissman2020}. Enquanto medidas nos processos fornecedores buscam reduzir a emissão por unidade produzida, mudanças de projeto podem reduzir desperdícios e necessidades de materiais para entregar o mesmo serviço.

A divulgação verificável das emissões dos materiais e a preferência por produtos de menor emissão nas compras públicas são instrumentos para estimular a demanda por alternativas de menor impacto \parencite[seções~6.6 e~6.7]{rissman2020}. O diagnóstico intersetorial identifica origens a considerar na avaliação da cadeia, mas não revela contratos ou poder de negociação. A especificação de materiais requer informações de fornecedores e avaliação de ciclo de vida; estimar o alcance das compras públicas requer conhecer a participação governamental na demanda da construção. O nó de administração pública não mede essa participação.

Para o setor de energia, o perfil de saída mostra que, mantidas a estrutura produtiva e a demanda do modelo, uma redução de sua intensidade emissora reduziria as emissões atribuídas a diversos produtos finais. Essa leitura dialoga com a articulação entre fornecimento energético de menor emissão e mudanças nos processos industriais \parencite[seções~4.2 e~7.2]{rissman2020}. Os dados agregados não permitem escolher tecnologias nem estimar a eficácia ou o custo de instrumentos. A distribuição normativa dos esforços de mitigação requer critérios adicionais à atribuição contábil.

As interpretações ambientais permanecem condicionadas à base utilizada. A interpolação, a cobertura não reconciliada e a diferença de preços limitam a aplicação dos contrastes setoriais à economia brasileira observada. Essas restrições alcançam os volumes e as comparações entre atividades, podendo também afetar os perfis ponderados e a seleção de relações. A consistência contábil e a sensibilidade a $\alpha$ verificam aspectos distintos da validade das emissões físicas. Antes de sustentar conclusões ambientais definitivas, é necessário reconciliar a fronteira, a alocação setorial e os denominadores monetários dos coeficientes ou substituí-los por uma extensão compatível com 2015.

\clearpage
\section{Conclusões}

No cenário analisado, as emissões de energia e água/resíduos atribuídas a outros produtos finais superam as emissões provenientes de outras origens e atribuídas aos produtos finais de energia e água/resíduos. Em construção ocorre o inverso, ampliando sua conta pela demanda em relação à produção. A identificação das contrapartes esclarece esses saldos e distingue perfis de volumes próximos com concentrações diferentes, apesar da conectividade semelhante. Construção também mostra a necessidade de separar as emissões atribuídas aos produtos finais da própria atividade das relações intersetoriais. A leitura conjunta das contas e dos perfis permite compreender a composição das diferenças entre produção e demanda.

O backbone distingue relações destacadas para a origem, para o destino ou para ambos. Os casos mostram que essa importância pode ser assimétrica e que a estabilidade dos papéis examinados convive com mudanças na seleção e na cobertura.

Na cadeia da construção, os resultados oferecem elementos para considerar conjuntamente as emissões nos processos fornecedores, a eficiência no uso de materiais e critérios de menor emissão nos projetos e nas aquisições. A interpretação sustenta a discussão dessas estratégias, sem estabelecer a eficácia ou o custo de instrumentos. A aplicação a Santa Catarina requer contas e perfis próprios, e conclusões empíricas definitivas sobre o Brasil dependem da validação ou substituição da extensão ambiental.

\clearpage
\begingroup
\setstretch{1}
\raggedright
\printbibliography[title={Referências}]
\endgroup
\input{anexos.tex}
\end{document}
